% ch04-security
\section{الأمان بعُمق (Defense in Depth)}

في الفصل السابق شرحنا كيف يعرف النظام \emph{مَن أنت} (التوثيق والرموز). في هذا
الفصل نجيب عن سؤال مختلف: ماذا يحدث \emph{بعد} أن نعرفك؟ كيف نحمي بياناتك وبيانات
غيرك حتى لو نجح مهاجم في اختراق طبقة واحدة؟ الجواب هو فلسفة اسمها \en{Defense in
Depth}، أو «الأمان بعُمق». وسنسير فيها \emph{واحدة واحدة}.

\begin{idea}[الفكرة: طبقات مستقلّة لا حائط واحد]
تخيّل قلعة قديمة: خندق، ثم سور، ثم بوابة، ثم حرّاس بالداخل. لو قفز أحدهم فوق
الخندق، ما زال السور يوقفه. الأمان بعُمق نفس الفكرة في البرمجيات: بدلًا من الاعتماد
على ضابط أمان \emph{واحد}، نضع عدّة ضوابط \textbf{مستقلّة}، كلٌّ منها يعمل وحده. فشل
طبقة لا يعني سقوط النظام، لأن الطبقة التالية ما زالت قائمة. هذا يحوّل «ثغرة واحدة =
كارثة» إلى «ثغرة واحدة = إزعاج محدود».
\end{idea}

\subsection{خريطة الطبقات كاملةً}

قبل الغوص في التفاصيل، هذه صورة عامّة لكل طبقات الحماية في \en{ScholarPath}. كل صفّ
في الجدول ضابط مستقلّ؛ وباقي الفصل يشرح الصفوف الحسّاسة منها بالتفصيل (التشفير،
الرفع، الفحص، الخصوصيّة، والتدقيق).

\begin{center}
\begin{tabular}{L{3.4cm} L{9.4cm}}
\toprule
\textbf{\color{spPrimary}الطبقة} & \textbf{\color{spPrimary}الضابط} \\
\midrule
النقل \en{(Transport)} & \en{HTTPS} في كل مكان؛ رمز \en{JWT bearer} على الـ \en{API}؛ محاور \en{SignalR} موثَّقة \\
الهويّة \en{(Identity)} & \en{ASP.NET Core Identity} — كلمات سرّ مُجزّأة \en{(hashed)}، قفل عند التخمين، تأكيد بريد \\
الرموز \en{(Tokens)} & رمز وصول \en{JWT} قصير العمر + رموز تحديث \en{(refresh)} دوّارة تُخزَّن \emph{مُجزّأة} \\
التخويل \en{(Authorization)} & أدوار (\en{Student}/\en{ScholarshipProvider}/\en{Consultant}/\en{Admin}) + ملكيّة لكلّ مورد \\
الأسرار والمفاتيح & مزوّدات \en{Azure Key Vault} (مفتاح تشفير الحقول، مفتاح توقيع \en{JWT}) \\
البيانات المخزَّنة & تشفير حقول \en{AES-256-GCM} على الأعمدة الحسّاسة \\
الرفع \en{(Uploads)} & فحص-قبل-التخزين بـ \en{ClamAV}، بسياسة \en{fail-closed} \\
الخصوصيّة/\en{PII} & تنقية النصّ قبل الذكاء الاصطناعي؛ مراجع فضفاضة؛ تصدير/حذف \en{GDPR} \\
التكامل \en{(Integrity)} & مفاتيح \en{idempotency}، فهارس فريدة مُرشَّحة، \code{RowVersion} \\
قابليّة التدقيق & سجلّ \code{AuditLog} إلحاقيّ فقط (قبل/بعد بصيغة \en{JSON}) \\
الحذف الناعم & مُرشِّحات استعلام عامّة تُخفي الصفوف المحذوفة من كلّ استعلام \\
\bottomrule
\end{tabular}
\end{center}

\begin{why}
لاحظ أنّ الطبقات \emph{لا تتكلّم مع بعضها}. التشفير لا يعرف شيئًا عن فحص الفيروسات،
والتدقيق لا يعرف شيئًا عن الأدوار. هذا الاستقلال مقصود: لو ارتبطت الطبقات ببعضها،
لأصبح خطأ في واحدة قادرًا على تعطيل الباقي. الاستقلال هو ما يمنح «العُمق» معناه.
\end{why}

\subsection{التشفير عند التخزين: AES-256-GCM}

بعض الأعمدة في قاعدة البيانات حسّاسة جدًّا — مثل نبذة المستخدم \code{Biography} أو
الملاحظات الشخصيّة على الطلب \code{PersonalNotes}. لو تسرّبت نسخة احتياطيّة من قاعدة
البيانات، لا نريد أن يقرأ أحدٌ هذه الحقول. الحلّ: نشفّرها \emph{قبل} أن تُكتَب، على
مستوى التطبيق نفسه.

\begin{idea}[الفكرة: تشفير مُصادَق (AEAD) بمغلّف ذاتيّ الوصف]
نستخدم خوارزمية \en{AES-256-GCM}. كلمة \en{GCM} مهمّة: إنها خوارزمية
\textbf{مُصادَقة} \en{(AEAD)} تعطيك \emph{السرّيّة والسلامة معًا} في خطوة واحدة.
المفتاح 32 بايت (256 بت)، ولكلّ عمليّة تشفير \en{nonce} عشوائيّ طازج طوله 12 بايت،
ووسم مُصادَقة \en{(tag)} طوله 16 بايت. القيمة المخزَّنة «مغلّف ذاتيّ الوصف»:
البادئة الحرفيّة \code{enc:v1:} متبوعةً بـ \en{Base64} للسلسلة
(\en{nonce} ثم \en{tag} ثم \en{ciphertext}). لأنّ الـ \en{nonce} عشوائيّ، فإنّ تشفير
النصّ نفسه مرّتين يعطي ناتجًا مختلفًا في كلّ مرّة.
\end{idea}

المنفذ \code{IFieldEncryptionService} يُنفّذه الصنف \code{AesGcmFieldEncryptionService}.
هنشوف إزاي بيبني المغلّف عند التشفير:

\begin{english}
\begin{lstlisting}[language={[Sharp]C}]
public const string EnvelopePrefix = "enc:v1:";
private const int NonceSizeBytes = 12;  // AES-GCM standard nonce
private const int TagSizeBytes   = 16;  // authentication tag

public string? Encrypt(string? plaintext)
{
    if (plaintext is null) return null;                 // null -> null
    var plaintextBytes = Encoding.UTF8.GetBytes(plaintext);

    var nonce = new byte[NonceSizeBytes];
    RandomNumberGenerator.Fill(nonce);                  // fresh random nonce

    var tag = new byte[TagSizeBytes];
    var ciphertext = new byte[plaintextBytes.Length];
    using (var aes = new AesGcm(_key, TagSizeBytes))
        aes.Encrypt(nonce, plaintextBytes, ciphertext, tag);

    // Envelope payload: nonce (12) + tag (16) + ciphertext
    var payload = new byte[NonceSizeBytes + TagSizeBytes + ciphertext.Length];
    // ... copy nonce, then tag, then ciphertext into payload ...
    return EnvelopePrefix + Convert.ToBase64String(payload);
}
\end{lstlisting}
\end{english}

نشرح أهمّ الأسطر:
\begin{itemize}
\item \code{if (plaintext is null) return null} — القيمة الفارغة تمرّ كما هي؛ عمود
  فارغ يبقى فارغًا ولا نضيّع تشفيرًا عليه.
\item \code{RandomNumberGenerator.Fill(nonce)} — هنا يُولَد الـ \en{nonce} العشوائيّ
  الجديد. هذا السطر تحديدًا هو سرّ أنّ نفس النصّ لا يعطي أبدًا نفس التشفير.
\item \code{aes.Encrypt(nonce, plaintextBytes, ciphertext, tag)} — خطوة واحدة تُنتج
  النصّ المشفَّر \emph{و} وسم المُصادَقة معًا.
\item السطر الأخير يجمّع (\en{nonce} + \en{tag} + \en{ciphertext}) ويضع أمامها
  البادئة \code{enc:v1:} — فتصير القيمة تصف نفسها بنفسها.
\end{itemize}

وعلى الجانب الآخر، عند فكّ التشفير، يُتحقَّق من الوسم \en{(tag)} أوّلًا: أيّ عبثٍ
ببِتّ واحد، أو اقتطاع للقيمة، أو مفتاح خاطئ — يجعل الدالّة \code{Decrypt}
\textbf{ترمي استثناءً} بدلًا من أن تُرجع قمامة. هذا هو معنى «مُصادَق».

\begin{keyp}[نقطة مهمة: versioned + طرح آمن]
المقطع \code{v1} في البادئة ليس زينة — إنه \textbf{رقم إصدار}. لو دوّرنا المفتاح
مستقبلًا، نُدخل \code{enc:v2:} للقيم الجديدة بينما نظلّ قادرين على فكّ القيم القديمة
\code{v1}. وأهمّ من ذلك: أيّ قيمة \emph{لا} تبدأ بالبادئة (نصّ قديم كُتب قبل تفعيل
الميزة) تُرجَع كما هي دون تعديل، ثم يُعاد تشفيرها تلقائيًّا عند أوّل كتابة لها. هذا
«الطرح الآمن» يعني تفعيل التشفير على نظام حيّ \emph{بلا توقّف ولا هجرة بيانات}.
\end{keyp}

والأجمل أنّ بقيّة الكود لا تشعر بشيء من هذا. نربط التشفير بالعمود عبر \en{value
converter} في \en{EF Core}، فيصير النصّ صريحًا في الذاكرة ومشفَّرًا في القاعدة:

\begin{english}
\begin{lstlisting}[language={[Sharp]C}]
public sealed class EncryptedStringConverter : ValueConverter<string?, string?>
{
    public EncryptedStringConverter(IFieldEncryptionService encryption)
        : base(
            plaintext => encryption.Encrypt(plaintext),  // to the database
            stored    => encryption.Decrypt(stored))     // from the database
    { }
}
\end{lstlisting}
\end{english}

بمجرّد إسناد هذا المحوّل لعمود، تظلّ معالِجات الأوامر والاستعلامات تقرأ وتكتب نصًّا
عاديًّا، ولا تعلم أبدًا أنّ العمود مشفَّر عند التخزين. طبقة واحدة، ومسؤوليّة واحدة.

\begin{secnote}
لماذا \en{GCM} وليس تشفيرًا عاديًّا؟ التشفير العاديّ يحمي السرّيّة لكنه \emph{لا}
يكشف العبث. \en{GCM} يعطيك الاثنين. وميزته على تشفير القاعدة الشفّاف \en{(TDE)} أنّ
المفتاح لا يعيش أبدًا في الكود ولا في ملفّات الإعداد — بل يأتي من \en{Key Vault} —
وحتى نسخة احتياطيّة مسروقة أو مدير قاعدة بيانات فضوليّ لن يرى إلّا نصًّا مشفَّرًا
لهذه الأعمدة، وهو ما لا يضمنه تشفير القاعدة وحده بمجرّد أن يُرجِع استعلامٌ صفًّا.
\end{secnote}

\subsubsection{إدارة المفاتيح}

المفتاح يأتي من المنفذ \code{IFieldEncryptionKeyProvider}، وله مُنفِّذان: واحد يقرأ
من \en{Key Vault} في الإنتاج، وآخر محلّيّ للتطوير. وعند الإقلاع نتحقّق أنّ طوله
\textbf{32 بايت بالضبط}، وإلّا يرفض النظام أن يبدأ أصلًا. وهناك منفذ مستقلّ آخر هو
\code{IJwtKeyProvider} لمفتاح توقيع \en{JWT}. وضع كلّ مفتاح خلف منفذه الخاصّ هو ما
يجعل \textbf{تدوير المفاتيح} عمليّة تشغيليّة بسيطة، لا إعادة نشرٍ للكود.

\begin{qa}
«ماذا لو تسرّب المفتاح نفسه؟» — عندها يفشل هذا الضابط، لكن تذكّر فلسفة العُمق:
التشفير طبقة واحدة فقط. ما زالت هناك طبقات الهويّة والتخويل والتدقيق. ثم إنّ المفتاح
لا يعيش في الكود ولا في \en{git}؛ يعيش في \en{Key Vault} بصلاحيّات محدودة، ويمكن
تدويره فورًا إلى \code{enc:v2:} دون أن نفقد القيم القديمة.
\end{qa}

\subsection{رفع الملفات: مسار مُعتِم ومزوّدان}

المستخدمون يرفعون مرفقات: مستندات الطلبات، ملفّات الترقية، مرفقات المنتدى،
وتسجيلات الجلسات. المنفذ \code{IBlobStorageService} يُنفّذه \code{FileStorageService}،
وله مزوّدان يُختاران من الإعداد \code{Storage:Provider}: \en{Local} (قرص، للتطوير)
أو \en{AzureBlob} (تخزين \en{Azure}).

هنشوف إزاي بيتخزّن الملفّ وبيرجّع مساره:

\begin{english}
\begin{lstlisting}[language={[Sharp]C}]
public async Task<string> UploadAsync(
    Stream content, string fileName, string contentType,
    string container, CancellationToken ct)
{
    // Unique key keeps two same-named uploads from colliding.
    var key = $"{Guid.NewGuid():N}/{fileName}";

    if (UseAzure)
    {
        var blob = ContainerClient(container).GetBlobClient(key);
        await blob.UploadAsync(content, /* ...ContentType... */, ct);
        return $"azure:{container}/{key}";     // opaque path
    }
    // ... Local provider writes under LocalRoot()/container/key ...
    return $"local:{container}/{key}";
}
\end{lstlisting}
\end{english}

نشرح أهمّ الأسطر:
\begin{itemize}
\item \code{var key = \$"\{Guid.NewGuid():N\}/\{fileName\}"} — يُسبَق اسم الملفّ
  بـ \en{GUID} عشوائيّ، فملفّان بنفس الاسم لن يتصادما أبدًا، ولا يستطيع مستخدم أن
  يخمّن مسار ملفّ آخر.
\item القيمة المُرجَعة \code{azure:container/key} أو \code{local:container/key} هي
  \textbf{مسار مُعتِم} \en{(opaque)} يحمل اسم المزوّد بداخله — فيعمل التنزيل والحذف
  للمزوّدَين معًا دون لبس.
\end{itemize}

وحين نُنشئ حاوية \en{Azure}، نُنشئها بـ \code{PublicAccessType.None} — أي \emph{غير
عامّة أبدًا}. لا يوجد رابط عامّ للملفّ؛ كلّ تنزيل يمرّ عبر الـ \en{API} الموثَّق الذي
يتحقّق من صلاحيّة الطالب أوّلًا.

\begin{note}
المسار المُعتِم فكرة صغيرة بأثر كبير: الكود الأعلى يتعامل مع «سلسلة» واحدة يخزّنها في
عمود، ولا يهتمّ إن كان الملفّ على قرصٍ محلّيّ أو في \en{Azure}. تبديل المزوّد يصبح
تغيير إعدادٍ واحد، لا تغيير كود.
\end{note}

\subsection{الفحص قبل التخزين: ClamAV بسياسة fail-closed}

الملفّ المرفوع قد يكون خبيثًا. لذلك نفحصه \emph{قبل} أن نخزّنه، \textbf{وهو ما زال
تيّارًا في الذاكرة}. المنفذ \code{IFileScanService} يُنفّذه \code{ClamAvFileScanService}،
الذي يرسل البايتات إلى خادم \en{ClamAV} عبر بروتوكول \en{INSTREAM}. الأحكام ثلاثة:
\textbf{Clean} (نظيف، نكمل)، و\textbf{Infected} (مُصاب، نرفض ونسجّل التوقيع)،
و\textbf{ScanUnavailable} (تعذّر الفحص، نرفض).

هنشوف قلب الدالّة: كيف يترجم ناتج الفاحص إلى حُكم، وكيف يتصرّف عند الفشل:

\begin{english}
\begin{lstlisting}[language={[Sharp]C}]
switch (result.Result)
{
    case ClamScanResults.Clean:
        return new FileScanResult(FileScanVerdict.Clean, null);

    case ClamScanResults.VirusDetected:
        var signature = result.InfectedFiles?.FirstOrDefault()?.VirusName
                        ?? "unknown";
        return new FileScanResult(FileScanVerdict.Infected, signature);

    default: // Error / Unknown -> treat as unavailable
        return new FileScanResult(FileScanVerdict.ScanUnavailable, ...);
}
// ... and if the daemon is unreachable:
catch (Exception ex)   // socket timeout, protocol error, no daemon
{
    // Fail closed: an unscanned file is never stored.
    return new FileScanResult(FileScanVerdict.ScanUnavailable, "...");
}
\end{lstlisting}
\end{english}

نشرح المنطق:
\begin{itemize}
\item \code{Clean} فقط هو الذي يمرّ ويُخزَّن؛ التيّار يُعاد لبدايته لتُكتَب نفس
  البايتات المفحوصة.
\item \code{VirusDetected} يُرجِع الحُكم \code{Infected} مع توقيع الفيروس ويُسجَّل
  تحذيرًا؛ الملفّ يُرفَض.
\item السطر \code{catch (Exception ex)} هو جوهرة الفصل: أيّ تعذّرٍ في الوصول للخادم
  أو مهلة أو خطأ بروتوكول يُترجَم إلى \code{ScanUnavailable}، ومسار الرفع يرفض هذا
  الحُكم.
\end{itemize}

\begin{secnote}
هذه هي سياسة \en{fail-closed} (الفشل إلى الإغلاق). القاعدة: عند الشكّ، \emph{ارفض}.
لو تعطّل الفاحص، لا نقول «لا بأس، خزّن الملفّ»؛ بل نرفضه. لأنّ البديل — تخزين ملفّ
\emph{غير مفحوص} — يفتح نافذة زمنيّة قد يصبح فيها ملفّ خبيث في متناول اليد. الفحص
أوّلًا على التيّار، والرفض عند العجز، يغلقان هذه النافذة تمامًا. (والميزة مُتحكَّم بها
عبر \code{FileScanning:Enabled}، وفي التطوير يُستخدَم \code{NoOpFileScanService}.)
\end{secnote}

\subsection{بيانات الأشخاص (PII)}

بيانات الأشخاص \en{(Personally Identifiable Information)} — الإيميل، رقم الكارت، رقم
التليفون — نتعامل معها بحذر خاصّ، عبر عدّة طبقات في مسار الذكاء الاصطناعي وفي نموذج
البيانات وفي دورة الحياة.

\begin{idea}[الفكرة: تنقية قبل التخزين وقبل المزوّد]
حين يكتب المستخدم رسالة للمساعد الذكيّ، تمرّ الرسالة أوّلًا على دالّة
\code{RedactPii()} \textbf{قبل} أن نخزّنها \emph{وقبل} أن نرسلها لأيّ مزوّد خارجيّ.
فالنصّ الخام لا يغادر العمليّة أبدًا؛ ما نخزّنه (\code{AiInteraction.PromptText})
وما نرسله إلى \en{OpenAI} هو النسخة \emph{المُنقّاة}.
\end{idea}

هنشوف الدالّة الحقيقيّة وترتيب أنماطها:

\begin{english}
\begin{lstlisting}[language={[Sharp]C}]
internal static string RedactPii(string input)
{
    if (string.IsNullOrWhiteSpace(input)) return string.Empty;
    // Order matters: email, then card, then the most permissive phone last.
    var s = EmailRegex().Replace(input, "[redacted-email]");
    s     = CreditCardRegex().Replace(s, "[redacted-card]");
    s     = PhoneRegex().Replace(s, "[redacted-phone]");
    return s;
}

[GeneratedRegex(@"[\w\.-]+@[\w\.-]+\.\w{2,}")]      // email
private static partial Regex EmailRegex();
[GeneratedRegex(@"\b(?:\d[ -]*?){13,19}\b")]        // 13-19 digit card
private static partial Regex CreditCardRegex();
[GeneratedRegex(@"(?:\+?\d[\d\s\-\.]{8,}\d)")]      // 10+ digit phone
private static partial Regex PhoneRegex();
\end{lstlisting}
\end{english}

لماذا هذا الترتيب بالذات؟
\begin{itemize}
\item الإيميلات أوّلًا لأنها تتمركز حول علامة \code{@}، فيسهل تمييزها.
\item ثم الكروت: تسلسل من 13 إلى 19 رقمًا.
\item ثم التليفونات \textbf{أخيرًا}، لأنّ نمطها هو \emph{الأوسع} (أيّ 10 أرقام أو
  أكثر). لو جاء أوّلًا، لابتلع أرقام الكروت قبل أن يميّزها نمط الكارت. الترتيب هنا
  ليس تفصيلًا تجميليًّا؛ إنه صحّة.
\end{itemize}

وهناك طبقات مكمّلة حول الـ \en{PII}:
\begin{itemize}
\item \textbf{تخصيص بلا بيانات تواصُل} — الكتلة التعريفيّة التي نحقنها في المحادثة
  لتخصيص الإجابات تحمل حقائق أكاديميّة فقط (الجنسيّة، المستوى، المجال، المعدّل، الدول
  المفضّلة)، ولا تحمل إيميلًا ولا تليفونًا أبدًا.
\item \textbf{ذاكرة \en{AI} لكلّ مستخدم} — عند إعادة تشغيل تاريخ المحادثة نُرشّحه
  للمستخدم صاحب الطلب فقط، فحتى لو تسرّب \code{SessionId} لا يمكن أن يتسرّب نصّ
  محادثة مستخدمٍ آخر.
\item \textbf{عيّنات تدقيق التنقية} — الكيان \code{AiRedactionAuditSample} يخزّن
  عيّنات شهريّة من الرسائل المُنقّاة (حتى 50 عيّنة من الشهر المنصرم) لمراجعة بشريّة —
  حتى نُثبت أنّ المُنقّي يعمل فعلًا.
\item \textbf{مراجع فضفاضة \en{(loose FK)}} — الجداول عالية الحجم/التدقيق/التحليلات
  (المدفوعات، الإشعارات، الأصوات، الدردشة، سجلّ التدقيق، تفاعلات \en{AI}) تحمل
  \en{Guid} المستخدم \emph{بلا قيد مفتاح أجنبيّ}، فيمكن حذف مستخدم أو إخفاء هويّته
  دون أن ينكسر التاريخ تتابعيًّا.
\item \textbf{\en{GDPR}} — الكيان \code{UserDataRequest} يقود التصدير والحذف عبر
  وظيفتَي \code{DataExportJob} و\code{DataDeleteJob}؛ والنصّ الحرّ الحسّاس مشفَّر عند
  التخزين كما رأينا.
\end{itemize}

\begin{secnote}
لاحظ الفلسفة المتكرّرة: نقلّل ما نحتفظ به من الأساس. التنقية تحذف الـ \en{PII} قبل
أن يُكتَب، والمراجع الفضفاضة تجعل حذف المستخدم ممكنًا دون تدمير السجلّات المرتبطة.
أقلّ بيانات حسّاسة مخزَّنة = سطح هجوم أصغر.
\end{secnote}

\subsection{سجلّ التدقيق (AuditLog): إلحاقيّ فقط}

آخر طبقة نتناولها هي \textbf{قابليّة التدقيق}. كلّ فعل تجاريّ حسّاس يُسجَّل في جدول
\code{AuditLog} بصورة \en{append-only} (إلحاقيّ فقط) — نضيف صفوفًا ولا نعدّل ولا نحذف.
الخدمة \code{AuditService} تكتب الصفّ، وتحفظ الحالة \emph{قبل} و\emph{بعد} بصيغة
\en{JSON}:

\begin{english}
\begin{lstlisting}[language={[Sharp]C}]
var entry = new AuditLog
{
    ActorUserId  = currentUser.UserId, // null for system actions
    Action       = action,
    TargetType   = targetType,         // polymorphic: any entity kind
    TargetId     = targetId,
    BeforeJson   = beforeJson,          // state before the change
    AfterJson    = afterJson,           // state after the change
    IpAddress    = Trim(currentUser.IpAddress, 64),
    CorrelationId= Trim(currentUser.CorrelationId, 128),
    OccurredAt   = clock.UtcNow,
};
db.AuditLogs.Add(entry);
\end{lstlisting}
\end{english}

نشرح الحقول المهمّة:
\begin{itemize}
\item \code{ActorUserId} — مَن فعلها؛ ويكون فارغًا \en{(null)} إذا كان الفاعل هو
  النظام نفسه (وظيفة خلفيّة مثلًا).
\item \code{TargetType} و\code{TargetId} — مرجع \en{polymorphic}: يستطيع الصفّ أن
  يشير إلى \emph{أيّ} نوع كيان، فسجلٌّ واحد يخدم كلّ الكيانات.
\item \code{BeforeJson} و\code{AfterJson} — لقطتا الحالة قبل التغيير وبعده، فنعرف
  بالضبط \emph{ماذا} تغيّر، لا مجرّد أنّ شيئًا تغيّر.
\end{itemize}

\begin{keyp}[نقطة مهمة: لماذا «إلحاقيّ فقط»؟]
سجلّ التدقيق لا قيمة له إن أمكن تزويره. جعله \en{append-only} يعني أنّ التاريخ لا
يُعاد كتابته: مَن فعل ماذا، ومتى، وما القيمة قبل وبعد — كلّها محفوظة إلى الأبد. وحين
يجتمع هذا مع الحذف الناعم (الصفوف لا تُدمَّر بل تُخفى)، تحصل على أثرٍ كامل قابل
للتتبّع لأيّ نزاع أو تحقيق لاحق.
\end{keyp}

\begin{summary}[الخلاصة في سطور]
\begin{itemize}
\item \textbf{الأمان بعُمق} = طبقات ضوابط \emph{مستقلّة}؛ فشل طبقة لا يُسقط النظام،
  لأنّ الطبقة التالية ما زالت قائمة.
\item \textbf{التشفير عند التخزين} بـ \en{AES-256-GCM}: مفتاح 32 بايت، \en{nonce}
  عشوائيّ 12 بايت، وسم 16 بايت، ومغلّف ذاتيّ الوصف \code{enc:v1:}؛ مُصادَق فأيّ عبثٍ
  يُفشِل فكّ التشفير، ومُصدَّر \en{(versioned)} فالطرح آمن بلا توقّف.
\item \textbf{رفع الملفّات} عبر منفذٍ بمزوّدَين ومسارٍ مُعتِم بادئته \en{GUID}،
  والحاويات \code{PublicAccessType.None} — لا رابط عامّ أبدًا.
\item \textbf{الفحص قبل التخزين} بـ \en{ClamAV} على التيّار، بسياسة \en{fail-closed}:
  عند الشكّ ارفض؛ لا نُخزّن ملفًّا غير مفحوص.
\item \textbf{الـ \en{PII}} تُنقَّى قبل التخزين وقبل المزوّد (بترتيب: إيميل ثم كارت
  ثم تليفون)، مع تخصيصٍ بلا بيانات تواصُل، ومراجع فضفاضة، وتصدير/حذف \en{GDPR}.
\item \textbf{سجلّ التدقيق} \code{AuditLog} إلحاقيّ فقط، يحفظ الحالة قبل/بعد
  بـ \en{JSON} — أثرٌ لا يُزوَّر لكلّ فعلٍ حسّاس.
\end{itemize}
\end{summary}
